\documentclass[a4paper, serif, 10pt]{moderncv}

%% moderncv
% style and color
\moderncvstyle{banking}
\moderncvcolor{black}

%% page layout/margins
\usepackage[top=2.5cm, bottom=2.5cm, left=1.5cm, right=1.5cm]{geometry}

\nopagenumbers{}

%% Babel config for French language
% ref:
% - https://latex3.github.io/babel/guides/locale-french.html
\usepackage[french]{babel}

%% fontspec
\usepackage{fontspec}

\IfFontExistsTF{Linux Libertine}{
  \setmainfont[Ligatures={TeX, Common}, Numbers=OldStyle]{Linux Libertine}
}{}
\IfFontExistsTF{Linux Libertine O}{
  \setmainfont[Ligatures={TeX, Common}, Numbers=OldStyle]{Linux Libertine O}
}{}

%% CTeX
% CJK support, used to show CJK characters in the resume
%
% - fontset=none: disable builtin fontset but instead set the CJK font manually
% - heading=false: disable ctex heading
% - punct=kaiming: use kaiming punctuations styles for CJK
% - scheme=plain: use plain scheme, do not override \normalsize font size
% - space=auto: space settings for CJK characters
%
% ref:
% - http://ctan.mirrorcatalogs.com/language/chinese/ctex/ctex.pdf
\usepackage[UTF8, heading=false, punct=kaiming, scheme=plain, space=auto]{ctex}

\IfFontExistsTF{Noto Serif CJK SC}{
  \setCJKmainfont{Noto Serif CJK SC}
}{}
\IfFontExistsTF{Noto Sans CJK SC}{
  \setCJKsansfont{Noto Sans CJK SC}
}{}

%% line spacing
\usepackage{setspace}
\setstretch{1.125}

%% URL styling - use normal text instead of monospace
\urlstyle{same}

% Auto-underline all links
% Defer until \begin{document} so hyperref is already loaded
\AtBeginDocument{%
  \let\oldhref\href
  \renewcommand{\href}[2]{\oldhref{#1}{\underline{#2}}}%
}

%% Patch \httplink and \httpslink to handle full URLs
% The original moderncv macros always prepend http:// or https:// to URLs.
% This causes issues when the URL already has a protocol (e.g., https://example.com)
% resulting in malformed URLs like https://https://example.com.
% This patch checks if the URL already contains "://" and uses it directly if so.
% Note: We use \str_set:Nx (x-expansion) to expand macros like \@homepage first.
\ExplSyntaxOn
\str_new:N \g_yamlresume_url_str

\RenewDocumentCommand{\httpslink}{O{}m}{
  \str_set:Nx \g_yamlresume_url_str {#2}
  \str_if_in:NnTF \g_yamlresume_url_str {://}
    { \url{#2} }
    { \url{https://#2} }
}

\RenewDocumentCommand{\httplink}{O{}m}{
  \str_set:Nx \g_yamlresume_url_str {#2}
  \str_if_in:NnTF \g_yamlresume_url_str {://}
    { \url{#2} }
    { \url{http://#2} }
}
\ExplSyntaxOff

\name{Camille Moreau}{}
\title{Designer UX senior — conception de produits numériques centrés utilisateur}
\phone[mobile]{+33 6 12 34 56 78}
\email{camille.moreau@email.fr}
\homepage{https://www.camillemoreau.fr}

\address{12 Rue de la Paix, Paris, Île-de-France, France, 75002}{}{}

\extrainfo{{\small \faLinkedin}\ \href{https://www.linkedin.com/in/camillemoreau}{@camillemoreau} {} {} {} • {} {} {}
{\small \faBehance}\ \href{https://www.behance.net/camillemoreau}{@camillemoreau} {} {} {} • {} {} {}
{\small \faDribbble}\ \href{https://dribbble.com/camillemoreau}{@camillemoreau}}

\begin{document}

\maketitle

\section{Informations de base}

\cvline{}{Designer UX senior avec plus de 7 ans d'expérience dans la conception de produits numériques pour les secteurs de la santé, du e-commerce et des services financiers. Je m'applique à comprendre les besoins réels des utilisateurs pour concevoir des interfaces intuitives, accessibles et désirables.
%
\begin{itemize}
\item Expertise en recherche utilisateur, prototypage interactif et design system
\item Approche centrée utilisateur, guidée par les données et les tests d'utilisabilité
\item Collaboration étroite avec les développeurs, les product managers et les métiers
\item Passionnée par l'accessibilité numérique et le design inclusif
\end{itemize}}
\section{Formation}

\cventry{sept. 2014–juil. 2016}
        {Master, Design d'interaction et expérience utilisateur, Score : 3.8}
        {L'École de design Nantes Atlantique}
        {\url{https://www.lecolededesign.com/}}
        {}
        {\begin{itemize}
\item Formation centrée sur la conception d'interfaces et de services numériques innovants
\item Projet de fin d'études sur la conception d'un parcours de soin connecté pour patients chroniques
\item Stage de fin d'études en agence UX, avec création de tests utilisateurs et de wireframes
\end{itemize}
\textbf{Cours} : Ergonomie cognitive, Design thinking, Prototypage interactif, Recherche utilisateur, Design visuel, Méthodes agiles}

\cventry{sept. 2011–juil. 2014}
        {Licence, Arts plastiques — Design}
        {Université de Rennes 2}
        {\url{https://www.univ-rennes2.fr/}}
        {}
        {\begin{itemize}
\item Acquisition des fondamentaux du design visuel, de l'histoire de l'art et de la sémiologie de l'image
\item Réalisation de projets graphiques individuels et collectifs, alliant créativité et rigueur
\end{itemize}}
\section{Expérience professionnelle}

\cventry{janv. 2022–Present}
        {Designer UX senior}
        {Cortex Design}
        {\url{https://www.cortex-design.fr}}
        {}
        {\begin{itemize}
\item Pilote la conception UX/UI de la plateforme de télémédecine pour un client du secteur de la santé
\item Anime des ateliers de co-design avec les équipes produit, technique et médicale
\item Développe et fait évoluer le design system multi-produits (Figma, Storybook)
\item Planifie et réalise des tests d'utilisabilité, puis formalise des recommandations d'amélioration
\item Accompagne l'équipe dans l'application des critères d'accessibilité RGAA et WCAG 2.1
\end{itemize}
\textbf{Mots-clés} : Recherche utilisateur, Tests d'utilisabilité, Design system, Accessibilité, Figma}

\cventry{mars 2019–déc. 2021}
        {Product Designer}
        {Numen Studio}
        {\url{https://www.numen-studio.fr}}
        {}
        {\begin{itemize}
\item Conception de parcours utilisateurs et de wireframes pour des applications web et mobiles
\item Réalisation de prototypes interactifs présentés lors de recettes avec les clients
\item Mise en place d'un design system interne, réduisant de 30 \% le temps de production
\item Participation à la refonte de l'application mobile d'une grande enseigne de distribution
\item Contribution aux tests utilisateurs en laboratoire et à distance
\end{itemize}
\textbf{Mots-clés} : Wireframing, Prototypage, User journey, Design system, Sketch}

\cventry{sept. 2016–févr. 2019}
        {UI/UX Designer}
        {Pixelea}
        {\url{https://www.pixelea.fr}}
        {}
        {\begin{itemize}
\item Conception d'interfaces responsives pour des sites vitrines et des outils intranet
\item Création d'identités visuelles et de chartes graphiques pour des PME
\item Collaboration avec les développeurs pour assurer la faisabilité technique des maquettes
\item Réalisation de supports de présentation pour les réponses aux appels d'offres
\end{itemize}
\textbf{Mots-clés} : UI design, Intégration HTML/CSS, Identité visuelle, Responsive design}
\section{Langues}

\cvline{Français}{Niveau natif ou bilingue}
\cvline{Anglais}{Niveau professionnel complet \hfill \textbf{Mots-clés} : TOEIC 950}
\section{Compétences}

\cvline{Recherche utilisateur}{Expert \hfill \textbf{Mots-clés} : Entretiens utilisateurs, Personas, Parcours utilisateur, Tests d'utilisabilité, Analyse de données}
\cvline{Prototypage et design d'interface}{Expert \hfill \textbf{Mots-clés} : Figma, Sketch, Adobe XD, InVision, Principle}
\cvline{Design system et accessibilité}{Avancé \hfill \textbf{Mots-clés} : Design tokens, Composants, WCAG 2.1, RGAA}
\cvline{Collaboration et méthodes agiles}{Avancé \hfill \textbf{Mots-clés} : Scrum, Kanban, Lean UX, Ateliers de co-design}
\section{Récompenses}

\cventry{oct. 2021}
        {UX Design Awards Europe}
        {Prix UX Design – Mention accessibilité}
        {}
        {}
        {Récompense le projet « Mon Suivi Santé » pour sa conception inclusive et son interface accessible aux utilisateurs seniors.}
\section{Certificats}

\cventry{mars 2020}
        {Nielsen Norman Group}
        {Certified User Experience Professional (CUP)}
        {\url{https://www.nngroup.com/}}
        {}
        {}

\cventry{août 2021}
        {Coursera}
        {Google UX Design Professional Certificate}
        {\url{https://www.coursera.org/}}
        {}
        {}
\section{Publications}

\cventry{juin 2022}
        {Améliorer l'accessibilité des interfaces mobiles pour les seniors}
        {Les Cahiers du Design}
        {\url{https://www.lescahiersdudesign.fr/}}
        {}
        {Article de synthèse présentant les bonnes pratiques d'accessibilité numérique et les retours d'expérience de tests utilisateurs menés auprès de personnes âgées.}
\section{Projets}

\cventry{sept. 2021–Present}
        {Conception UX/UI d'une application mobile de suivi médical pour patients chroniques.}
        {Mon Suivi Santé}
        {\url{https://www.monsuivisante.fr}}
        {}
        {\begin{itemize}
\item Conduite d'entretiens utilisateurs et de tests d'utilisabilité
\item Création de l'architecture de l'information et des parcours clés
\item Conception d'une interface accessible conforme aux normes WCAG 2.1
\end{itemize}
\textbf{Mots-clés} : Mobile, Santé, Accessibilité, Figma, Tests utilisateurs}

\cventry{janv. 2023–juin 2023}
        {Refonte de l'expérience d'achat d'une marque de décoration éco-responsable.}
        {Maison Verte}
        {\url{https://www.maisonverte.fr}}
        {}
        {\begin{itemize}
\item Audit UX et analyse des parcours d'achat existants
\item Conception d'une interface épurée mettant en avant les engagements éco-responsables
\item Mise en place d'un design system et de composants réutilisables
\end{itemize}
\textbf{Mots-clés} : E-commerce, UX/UI, Design system, Conversion}
\section{Centres d'intérêt}

\cvline{Design social}{Design inclusif, Accessibilité numérique, Éco-conception}
\cvline{Arts et culture}{Photographie, Musées, Littérature}
\section{Bénévolat}

\cventry{sept. 2020–Present}
        {Bénévole UX}
        {Designers Éthiques}
        {\url{https://www.designersethiques.org/}}
        {}
        {Accompagnement d'associations à but non lucratif dans l'amélioration de leurs outils numériques et la formation de leurs équipes aux bonnes pratiques d'accessibilité.}

\end{document}